% Subsection 4.4.3: CoT 对照实验设计

本节与前后文的衔接关系如下：
\begin{itemize}
    \item \textbf{4.4.1--4.4.2}：已给出 CoT-ZS / CoT-FS 的流程与选池；
    \item \textbf{4.1 节（1a）（1b）}：集中给出各组评测时模型见到的 messages 字面骨架（含 Direct 与 CoT 训练侧对照），便于将 H1--H10 与「输入长什么样」逐组对应；
    \item \textbf{本节任务}：将 H2、H3、H5、H6、H8、H10 共 6 组 CoT 评测按设问重组为三块对照，作为第五章归因的直接依据。
\end{itemize}

\subsubsection{推理侧主效应（训练侧固定为 Base）}
以 H1 为共同基线，抽取不依赖任何 LoRA 微调的 CoT 推理侧贡献：
\begin{itemize}
    \item \textbf{$H2 - H1$}：CoT-ZS 相对 Direct 的主效应，回答「零示例推理是否在基座上就有增益」；
    \item \textbf{$H3 - H1$}：CoT-FS($k$) 相对 Direct 的主效应，回答「加 $k$ 条示例是否进一步改进基座」；
    \item \textbf{$H3 - H2$}：few-shot 相对 zero-shot 的边际贡献，隔离示例本身的增益；
    \item \textbf{对照解释}：三条差分共享同一训练侧（Base），仅推理侧变化，故组内显著差异可归因于 prompt 结构。
\end{itemize}

\subsubsection{训练 $\times$ 推理交互效应}
固定 LoRA 作为训练侧、CoT 作为推理侧，通过差分差分量（difference-in-differences）判断两者是否正交可加：
\begin{itemize}
    \item \textbf{$(H8 - H7) - (H2 - H1)$}：在 \texttt{LoRA\_cot\_zs} 与 Base 上「加 CoT-ZS 推理」两条边际是否一致；
    \item \textbf{$(H10 - H9) - (H3 - H1)$}：在 \texttt{LoRA\_cot\_fs($k$)} 与 Base 上「加 CoT-FS($k$) 推理」两条边际是否一致（$k$ 匹配约束下）；
    \item \textbf{$(H5 - H4) - (H2 - H1)$}：在 \texttt{LoRA\_code} 与 Base 上「加 CoT-ZS 推理」两条边际是否一致，额外控制了「未监督 CoT 格式」的训练侧；
    \item \textbf{$(H6 - H4) - (H3 - H1)$}：同上，推理侧换为 CoT-FS($k$)；
    \item \textbf{对照解释}：直接对应 4.1 节「LoRA 与 CoT 是否正交可加」的设问，符号与量级即析因设计的主落点。
\end{itemize}

\subsubsection{Few-shot 边际在训练侧上的一致性}
只变化训练侧，观察 few-shot 的边际贡献是否稳定：
\begin{itemize}
    \item \textbf{$H3 - H2$（Base）vs $H6 - H5$（LoRA\_code）}：两条「$k$ 例相对 $0$ 例」的边际在\textbf{未显式监督 CoT 格式}的训练侧上是否一致；
    \item \textbf{$H10 - H8$（匹配 $k$ 的 LoRA\_cot）}：训练与推理 $k$ 同步放大时，CoT 的净收益相对纯 CoT-ZS 路径的提升量；
    \item \textbf{设问}：few-shot 是否为与训练侧无关的「结构性增益」；
    \item \textbf{读法}：若两条边际量级接近，则 few-shot 主要作用于推理侧；若差异显著，则训练侧调节其有效性，结论中需分情况讨论。
\end{itemize}

\subsubsection{对照汇总}
\begin{itemize}
    \item 三组设计均为「单因子变动、其余固定」，覆盖全部与 CoT 相关的显著差分读法；
    \item 结果文件命名遵循表~\ref{tab:ch4-matrix} 末列；
    \item 第五章仅需按表取数作差，无需补做实验。
\end{itemize}